\documentclass[10pt,a4paper]{article}
\usepackage[spanish,es-nodecimaldot]{babel}
\usepackage[utf8]{inputenc}
\usepackage[T1]{fontenc}
\usepackage{lmodern}
\usepackage[top=1.1cm,bottom=1.1cm,left=1.5cm,right=1.5cm]{geometry}
\usepackage{amsmath,amssymb}
\usepackage[table]{xcolor}
\usepackage{booktabs,array,multirow}
\usepackage{enumitem}
\usepackage[normalem]{ulem}
\usepackage{pgfplots}
\pgfplotsset{compat=1.18}

\setlength{\parindent}{0pt}
\setlength{\parskip}{2pt}
\pagestyle{empty}

\AtBeginDocument{%
  \setlength{\abovedisplayskip}{3pt}%
  \setlength{\belowdisplayskip}{3pt}%
  \setlength{\abovedisplayshortskip}{2pt}%
  \setlength{\belowdisplayshortskip}{2pt}%
}

\begin{document}

{\large\bfseries\uline{Hipotesis relativa a varias proporciones}\par}
\vspace{2pt}

Se prueba que todas las proporciones de distintas poblaciones binomiales sean significativamente iguales:

\hspace{1.2cm}\textbf{H0:} $p_1 = p_2 = \ldots = p_k = p$ \\
\hspace*{1.2cm}\textbf{H1:} que al menos 1 no lo sea

\textbf{Método 1:} \ Si \textbf{n.p} y \textbf{(1-p).n} mayores a \textbf{5}: \\
La D.Binomial se aproxima a D.Normal, entonces, \uline{se calcula para cada población}:
\[
z_i = \frac{x_i - n_i\cdot p_i}{\sqrt{n_i\cdot p_i\cdot\left(1-p_i\right)}}
\]

\begin{minipage}{0.70\textwidth}
- El cuadrado de una v. aleatoria con dist. normal, es una v. aleatoria con dist. \textbf{chi-cuadrado} con 1 gl \\
- La suma de \textbf{k} v aleatorias con distribuciones \textbf{chi-cuadrado} con 1gl dan otra v. aleatoria \textbf{chi-cuadrado} con \textbf{k} gl. \hspace{0.8cm} gl=grados de libertad
\end{minipage}\hfill
\begin{minipage}{0.25\textwidth}
\[
Z_i^2 \sim \chi_1^2
\]
\end{minipage}

Dado que las $p_i$ son todas iguales por hipótesis nula y $p_n = x_n/n_n$, se pueden sustituir por:
\[
\widehat{p} = \frac{x_1 + x_2 + \cdots + x_k}{n_1 + n_2 + \cdots + n_k}
\qquad \text{suma de proporciones}
\]

$\bullet$ Obteniendo el estadístico:

\begin{minipage}{0.44\textwidth}
\[
\chi^2 = \sum_{i=1}^{k} \frac{\left(x_i - n_i\cdot\widehat{p}\,\right)^2}{n_i\cdot\widehat{p}\cdot\left(1-\widehat{p}\,\right)}
\]
\end{minipage}
\begin{minipage}{0.26\textwidth}
{\small Dif grandes se mueve hacia la der, dif chicas hacia la izq.\par
\textbf{v=k-1} grados de libertad\par}
\end{minipage}
\begin{minipage}{0.29\textwidth}
{\small fracasos esperados $n_i(1-\widehat{p}\,)$\par
éxitos esperados $n_i\widehat{p}$\par}
\end{minipage}

\begin{center}
\vspace{-6pt}
\begin{tikzpicture}
\begin{axis}[
  width=6.4cm, height=3.0cm,
  axis lines=left, axis line style={-},
  xmin=0, xmax=15, ymin=0, ymax=0.21,
  xtick=\empty, ytick=\empty,
  clip=false
]
\addplot[domain=0.05:14.5, samples=120, thick, blue] {x*exp(-x/2)/4};
\draw[thick] (axis cs:8.4,0) -- (axis cs:8.4,0.09);
\node[above, font=\scriptsize] at (axis cs:8.4,0.09) {$\chi^2_{\alpha}$};
\node[font=\tiny] at (axis cs:9.1,0.012) {$\alpha$};
\node[align=left, font=\tiny, text=blue!60!violet] at (axis cs:4.4,0.035) {valores chicos\\ caen acá};
\node[align=left, font=\tiny, text=violet] at (axis cs:12.3,0.065) {valores grandes\\ caen acá};
\end{axis}
\end{tikzpicture}
\vspace{-8pt}
\end{center}

\begin{center}
\vspace{-4pt}
\begin{tabular}{l r r r r}
\toprule
\rowcolor{cyan!10} \textbf{Resultado} & \textbf{A} & \textbf{B} & \textbf{C} & \textbf{Total} \\
Desmoronado & 41 & 27 & 22 & 90 \\
Resto intacto & 79 & 53 & 78 & 210 \\
\midrule
Total & 120 & 80 & 100 & 300 \\
\bottomrule
\end{tabular}
\end{center}

Se quiere determinar si la probabilidad de desmoronamiento es la misma para los materiales \textbf{A}, \textbf{B} y \textbf{C}. En cada material hay dos resultados posibles: desmoronado o resto intacto.

\textbf{metodo 1}
\[
\chi^2 = \sum_{i=1}^{3} \frac{\left(x_i - n_i\widehat{p}\,\right)^2}{n_i\widehat{p}\,(1-\widehat{p}\,)}
\qquad
\chi^2 = \frac{(x_1-n_1p)^2}{n_1p(1-p)} + \frac{(x_2-n_2p)^2}{n_2p(1-p)} + \frac{(x_3-n_3p)^2}{n_3p(1-p)}
\qquad
\nu = k-1 = 3-1 = 2.
\]

\textbf{Método 2:} para simplificar el cálculo de chi cuadrado se hace una tabla

\begin{center}
\vspace{-4pt}
\begin{tabular}{l | c | c | c | c | l}
 & Muestra 1 & Muestra 2 & $\cdots$ & Muestra k & Total \\
\cline{2-5}
éxitos & $x_1$ & $x_2$ & $\cdots$ & $x_k$ & $x$ \\
\cline{2-5}
fallas & $n_1 - x_1$ & $n_2 - x_2$ & $\cdots$ & $n_k - x_k$ & $n-x$ \\
\cline{2-5}
Total & $n_1$ & $n_2$ & $\cdots$ & $n_k$ & $n$ \\
\end{tabular}
\vspace{-6pt}
\end{center}

$\bullet$ Cada celda es una frecuencia observada.

\begin{minipage}{0.55\textwidth}
\textbf{Frecuencias esperadas}
\[
\text{Éxitos } \ \frac{n_j\cdot x}{n}
\qquad\quad
\text{Fracasos } \ \frac{n_j\cdot (n-x)}{n}
\]
\[
\chi^2 = \sum_{i=1}^{2}\sum_{j=1}^{k} \frac{\left(o_{i,j} - e_{i,j}\right)^2}{e_{i,j}}
\qquad \textbf{v=k-1} \text{ GL}
\]
\end{minipage}\hfill
\begin{minipage}{0.40\textwidth}
{\small
$n_j$ = total columna j \\
$x$ = total de la fila exitos \\
$n-x$ = total fila fallas \\
$n$ = total final
}
\end{minipage}

$\bullet$ Este método es lo mismo que aplicar la fórmula del chi cuadrado que se vió más arriba.

\textbf{Este método organiza los mismos datos en una tabla $2\times 3$: dos resultados y tres materiales} \hfill $\nu = k-1 = 3-1 = 2.$

\textbf{metodo 2}
\[
\chi^2 = \sum_{i=1}^{2}\sum_{j=1}^{3} \frac{\left(O_{ij} - E_{ij}\right)^2}{E_{ij}},
\qquad
E_{ij} = \frac{(\text{total fila } i)(\text{total columna } j)}{\text{total general}}.
\]

\begin{center}
\vspace{-4pt}
\begin{tabular}{l r r r r}
\toprule
\rowcolor{cyan!10} \textbf{Esperadas} & \textbf{A} & \textbf{B} & \textbf{C} & \textbf{Total} \\
Desmoronado & 36 & 24 & 30 & 90 \\
Intacto & 84 & 56 & 70 & 210 \\
\midrule
Total & 120 & 80 & 100 & 300 \\
\bottomrule
\end{tabular}
\end{center}

\end{document}
